% homework/hw-style.tex
% Shared house style for every Year 7 homework packet. \input this from the
% packet file, right after \documentclass and the answer-key switch.
%
% Do not fork this file per packet. If a packet needs a different look, the
% question is whether every packet should have it.
%
% The choices here are deliberate and were arrived at by printing and looking:
%   * No solid dark banners. Every header is a light tint with a coloured spine
%     on the left, so colour reads clearly without flooding a school printer.
%   * No marks and no mark totals. These are practice, not exam papers.
%   * Lato at 12pt with open leading, plus mathastext so inline maths is Lato
%     too. Without mathastext every "-6 + 4" is Computer Modern serif and the
%     sheet looks half-typeset.
%   * Writing lines are spaced for Year 7 handwriting, not adult margin notes.
%
% Provides: \wlines \blank \tickbox \sep \qmk-free marks (none) \hwsection
% \qhead \expr, the workspace box, and the remember / wordhelp / activity /
% yourturn coloured boxes. Palette matches the lesson decks exactly.

% -- Packages ----------------------------------------------------------------
\usepackage[T1]{fontenc}
\usepackage[utf8]{inputenc}
\usepackage[default]{lato}
\usepackage[a4paper,top=1.7cm,bottom=1.7cm,left=2.0cm,right=2.0cm,headheight=15pt]{geometry}
\usepackage{amsmath,amssymb}
\usepackage{xcolor}
\usepackage[most]{tcolorbox}
\usepackage{enumitem}
\usepackage{tikz}
\usetikzlibrary{arrows.meta}
\usepackage{array}
\usepackage{tabularx}
\usepackage{fancyhdr}
\usepackage{multicol}
\usepackage{needspace}
\usepackage{graphicx}

% Make maths use Lato too. Without this, every "-6 + 4" on the page is set in
% Computer Modern serif and the sheet looks half-typeset. Must come last.
\usepackage[italic]{mathastext}

\linespread{1.06}\selectfont

% -- The Learner's Book palette (same hex values as the lesson decks) ---------
\definecolor{cbteal}{HTML}{0087A8}
\definecolor{cbpurple}{HTML}{5C2483}
\definecolor{cborange}{HTML}{C25E12}
\definecolor{cbgreen}{HTML}{4A8B23}
\definecolor{cbred}{HTML}{C8102E}
\definecolor{cbblue}{HTML}{1A5FA8}
\definecolor{rulegray}{HTML}{9AA5AD}

% -- Answer lines and blanks -------------------------------------------------
% \wlines{n} draws n ruled writing lines, spaced wide enough for Year 7
% handwriting rather than for an adult with a fine pen.
\makeatletter
\newcommand{\wlines}[1]{%
  \par\vspace{0.75em}%
  \@tempcnta=0\relax
  \@whilenum\@tempcnta<#1\do{%
    \hbox to \linewidth{\color{rulegray}\leaders\hrule height 0.5pt\hfill}%
    \vspace{1.5em}%
    \advance\@tempcnta by 1\relax}%
  \par\vspace{-0.8em}%
}
\makeatother

% \blank[width] is an inline gap to write a short answer in.
\newcommand{\blank}[1][2.6cm]{\underline{\hspace{#1}}}

% A boxed space for showing working.
\newtcolorbox{workspace}[1][3.4cm]{%
  colback=white, colframe=rulegray, boxrule=0.5pt, arc=5pt,
  height=#1, left=6pt, right=6pt, top=4pt, bottom=4pt,
  before skip=9pt, after skip=11pt}

% An empty tick box.
\newcommand{\tickbox}{\fbox{\rule{0pt}{1.15em}\rule{1.15em}{0pt}}}

% A middot separator.
\newcommand{\sep}{\,\textperiodcentered\,}

% -- Section and question headings -------------------------------------------
% Light tint plus a fat coloured spine on the left. No solid dark banner: the
% colour reads clearly but barely costs the printer anything.
% needspace stops a heading stranding alone at the foot of a page. Kept modest:
% too large and it pushes headings over, leaving a worse hole than it prevents.
\newcommand{\hwsection}[2][cbteal]{%
  \needspace{8\baselineskip}%
  \par\vspace{13pt}%
  \noindent
  \begin{tcolorbox}[enhanced, nobeforeafter, width=\textwidth,
      colback=#1!7, colframe=#1!7, arc=5pt,
      borderline west={5pt}{0pt}{#1},
      left=14pt, right=10pt, top=7pt, bottom=7pt]
    {\large\bfseries\color{#1}#2}
  \end{tcolorbox}%
  \par\vspace{9pt}}

% \qhead[colour]{A2}{Moving along the line} -- a question heading.
\newcommand{\qhead}[3][cbteal]{%
  \needspace{6\baselineskip}%
  \par\vspace{11pt}%
  \noindent{\large\bfseries\color{#1}#2}\quad{\large\bfseries #3}\par\vspace{4pt}}

% -- Coloured boxes ----------------------------------------------------------
% Shared look: rounded, light fill, coloured rule, coloured title text on a
% pale title strip rather than a dark one.
\tcbset{hwbox/.style={
  enhanced, breakable, arc=6pt, boxrule=1pt,
  left=11pt, right=11pt, top=8pt, bottom=8pt,
  fonttitle=\bfseries, attach boxed title to top left={xshift=12pt, yshift=-8pt},
  boxed title style={arc=4pt, boxrule=1pt, left=7pt, right=7pt, top=2pt, bottom=2pt},
  top=13pt}}

% Orange: things to remember. Matches the deck's "Write This Down".
\newtcolorbox{remember}[1][Remember from the lesson]{hwbox,
  colback=cborange!6, colframe=cborange, coltitle=cborange,
  colbacktitle=white, title={#1}}

% Blue: English help. The barrier in this class is the wording, not the maths.
\newtcolorbox{wordhelp}[1][Word help]{hwbox,
  colback=cbblue!6, colframe=cbblue, coltitle=cbblue,
  colbacktitle=white, title={#1}}

% Purple: an activity or instruction.
\newtcolorbox{activity}[1][How to do this homework]{hwbox,
  colback=cbpurple!6, colframe=cbpurple, coltitle=cbpurple,
  colbacktitle=white, title={#1}}

% Green: your turn / a friendly nudge.
\newtcolorbox{yourturn}[1][Your turn]{hwbox,
  colback=cbgreen!7, colframe=cbgreen, coltitle=cbgreen!75!black,
  colbacktitle=white, title={#1}}

% -- Shorthands --------------------------------------------------------------
\newcommand{\dC}{\,^{\circ}\mathrm{C}}          % use inside maths: $-8\dC$
\newcommand{\key}[1]{{\color{cborange}\textbf{#1}}}
\newcommand{\eng}[1]{{\color{cbblue}\textbf{#1}}}

\setlist[enumerate]{itemsep=8pt, topsep=5pt, leftmargin=*}
\setlist[itemize]{itemsep=4pt, topsep=4pt, leftmargin=*}
\renewcommand{\arraystretch}{1.45}

% -- An expression the student is asked to work from -------------------------
\newcommand{\expr}[1]{%
  \tcbox[on line, colback=cbgreen!10, colframe=cbgreen!55, boxrule=1pt, arc=4pt,
         left=9pt, right=9pt, top=4pt, bottom=4pt]{\large\bfseries $#1$}}

% -- Running head and foot ---------------------------------------------------
% The packet overrides \fancyhead[L] with its own title.
\pagestyle{fancy}
\fancyhf{}
\fancyhead[L]{\small\color{black!55}Year 7\sep Homework}
\fancyhead[R]{\small\color{black!55}Name: \blank[3.4cm]}
\fancyfoot[C]{\small\color{black!55}Page \thepage}
\renewcommand{\headrulewidth}{0pt}
\renewcommand{\headrule}{\vspace{2pt}\hbox to\headwidth{\color{cbteal!45}\leaders\hrule height 1.2pt\hfill}}

